\section{State of the art and what \proj{} adds}
\label{sec:sota}

\subsection{Privacy-preserving linkage}

The dominant deployed approach encodes each identifier as a Bloom filter and compares filters by
Dice coefficient~\cite{schnell2009bloom}. It is fast, it needs no interaction, and it is used by
national statistical offices. It is also cryptanalysable. Attacks exploit the fact that the
encoding preserves frequency structure: common bigrams produce common bit patterns, and a graph
built on encoding similarity aligns with a graph built on plaintext similarity from any public name
frequency table~\cite{christen2018attack,haugsdal2022graphattack}. Reported recovery rates exceed
50 percent of surnames on realistic populations.

The field's response has been hardening: salting, bit balancing, record-level filters. Each raises
the cost of the attack without changing its structure, and none produces a statement of the form
"an adversary with this much side information learns at most this much". \proj{} takes the position
that the guarantee must be definitional rather than empirical, and that the cost of the guarantee
must be paid in a currency the epidemiologist can also see, namely the match quality.

\subsection{Linkage error in analysis}

Fellegi and Sunter gave the probabilistic framework in 1969~\cite{fellegi1969theory}, and the
statistical literature on estimation from linked files is small but sharp, beginning with the
regression corrections of Lahiri and Larsen~\cite{lahiri2005regression}. The applied literature
mostly ignores it. Registry studies typically report a match rate and then analyse the linked file
as observed data.

Our own preliminary work quantified what that costs~\cite{okparalange2025linkageerror}. Linkage
error is not classical measurement error, because the probability that a record is mislinked
depends on the same demographic variables that appear in the analysis model. Migrants, people who
change names, and people in large households are systematically harder to link. The bias therefore
concentrates in exactly the subgroups that the study is usually about.

\subsection{The gap}

No existing system reports, to the analyst, the quantities needed to bound that bias, because no
existing system knows them. The encoding is chosen for privacy, the threshold is chosen by a
custodian, and the resulting error rates are estimated, if at all, on a clerical review sample that
the analyst never sees. \proj{} closes the loop: the protocol that protects the identifiers is the
same protocol that certifies its own error rates, and the estimator consumes that certificate.

\subsection{Beyond the state of the art}

\begin{enumerate}
  \item A privacy definition for approximate string comparison that composes across a multi-party
    linkage, rather than a per-encoding heuristic.
  \item Sharp, computable bounds on linkage-induced bias for survival and causal estimands, not
    only for linear regression.
  \item A working three-custodian deployment, which is the point at which governance stops being a
    thought experiment~\cite{vestergaardnnamdi2023governance}.
\end{enumerate}
